% !TeX root = ../../Book2.tex
\OptionalChapter{多线性振荡积分与多线性 Fourier 限制}{b2:ch:25}
\BookChapterNumber{b2:ch:25}
\begin{chapterabstract}
本选读章研究相位与多输入投影的关系、法向横截及多线性增益。完整证明 Loomis–Whitney 不等式、平直横截延拓与分离抛物线片的双线性 \(L^2\) 公式；一般曲面多线性限制与解耦按正式来源标注为深层结构输入。
\end{chapterabstract}

\begin{learninggoals}\begin{itemize}
\item\label{b2:item:25-goal-1}能用可吸收相位构造排除虚假的统一参数衰减。
\item\label{b2:item:25-goal-2}能核对横截 determinant，并理解它与曲率的不同作用。
\item\label{b2:item:25-goal-3}能由频率和映射 Jacobi 行列式证明二维双线性模型。
\item\label{b2:item:25-goal-4}能准确读出多线性限制、波包及解耦中的尺度与损失。
\end{itemize}\end{learninggoals}

线性限制把许多频率方向混在一个输出里；多线性理论先保留不同方向的乘积结构。横截法向使几个相干管只在较小区域相遇，产生单输入估计看不见的增益。本章先建立可以完整计算的模型，再介绍其一般化所需的深层工具。

% !TeX root = ../../Book2.tex
\section{多线性积分形式}
\label{b2:sec:2501}

多线性估计保留多个输入之间的相互作用。当几个频率片的法向彼此分离时，乘积可能比单个延拓函数有更好的可积性；若相位可分别吸收到各输入中，则再大的参数也没有统一振荡增益。

\subsection{多线性型}
\label{b2:subsec:250101}

一个基本模型为
\[
 \Lambda_\lambda(f_1,\ldots,f_k)
 =\int_{\mathbb R^d}e^{2\pi\mathrm i\lambda P(x)}a(x)
                 \prod_{j=1}^k f_j(\pi_jx)\,\mathrm dx,
\]
其中 \(\pi_j\) 是固定线性投影，\(a\) 是紧支撑光滑振幅。讨论前应先用 Hölder 或具体投影几何确认积分有意义，再问是否有 \(\lambda^{-\epsilon}\) 的增益。
输出是多线性标量型，与\chapref{b2:ch:21}一个输入一个输出的算子不同。改变投影维数，会改变可允许的输入指数。

\ProseParagraphBreak
若投影相同，所有输入可先相乘，问题退化为单输入积分；若投影互补，输入的不同方向会限制相位抵消。多线性估计因而不仅依赖相位 Hesse 矩阵，还依赖它与整个投影族的关系。

\subsection{相位相互作用}
\label{b2:subsec:250102}

若相位可以写成 \(P=\sum_jP_j\circ\pi_j\)，则取
\(f_j(y)=e^{-2\pi\mathrm i\lambda P_j(y)}g_j(y)\) 就完全消去振荡，且各输入范数不变。若某些固定 \(g_j\) 使非振荡积分非零，就不可能有对所有输入一致的正幂次参数衰减。这称为相对于投影族的相位退化。

\ProseParagraphBreak
例如 \(\pi_1(x,y)=x\)、\(\pi_2(x,y)=y\)，相位 \(x^2+y^2\) 在这个意义下退化，虽然作为二维标量相位的 Hesse 矩阵 可逆。相位 \(xy\) 则不能写成纯 \(x\) 项加纯 \(y\) 项；其混合导数非零，并与\chapref{b2:ch:21}的 Fourier 型算子相联系。这个例子直接区分标量非退化与多线性非退化。

\subsection{缩放}
\label{b2:subsec:250103}

对无振幅的全空间齐次模型，令 \(f_j(y)\) 缩放为 \(f_j(ry)\)。左侧积分产生 \(r^{-d}\)，输入范数的乘积产生
\(r^{-\sum_j d_j/p_j}\)。尺度不变界因此要求
\(\sum_jd_j/p_j=d\)，其中 \(d_j\) 是第 \(j\) 个输入空间维数。
有固定紧支撑振幅时，尺度条件未必作为必要等式出现，但局部放大仍能检验极端集中。

\ProseParagraphBreak
抛物面还具有非各向同性缩放。切向频率缩小 \(\delta\)，法向频率变化 \(\delta^2\)，空间变量分别放大 \(\delta^{-1}\) 与 \(\delta^{-2}\)。多线性管束与解耦中的尺度都来自这一关系；在频率帽与空间球之间来回换算时，应始终保留平方关系。

\subsection{近似正交性}
\label{b2:subsec:250104}

若函数 \(u_j\) 的 Fourier 支撑有限重叠，则 Plancherel 给出
\(\|\sum_j u_j\|_2^2\le L\sum_j\|u_j\|_2^2\)。
对乘积 \(u_jv_k\)，Fourier 支撑位于两支撑的 Minkowski 和；即使原支撑互不相交，这些和仍可能重叠。因此多线性近似正交应核对输出支撑的重叠数，不能只核对各输入支撑。

\ProseParagraphBreak
在局部空间球上，乘光滑截断又会使 Fourier 支撑稍微膨胀。于是通常使用快速衰减权与有限扩大后的重叠数，而不是严格的全空间支撑等式。随机符号与平方函数提供另一种组合工具，但也要保持相关的算子常数一致。

% !TeX root = ../../Book2.tex
\section{横截性}
\label{b2:sec:2502}

横截性把方向分离量写成一个可核对的下界。最完整的横截模型是 \(d\) 个法向张成环境空间，此时平直曲面已经有非平凡的乘积估计，而不需要曲率。

\subsection{横截性}
\label{b2:subsec:250201}

设 \(S_j\) 为 \(d\) 个超曲面片，单位法向为 \(n_j(\xi_j)\)。一致横截条件是
\[
 |\det(n_1(\xi_1),\ldots,n_d(\xi_d))|\ge\nu>0
\]
对各片的任意点成立。它比“法向两两不同”强；在三维，三个互不平行的向量仍可能共面。下界 \(\nu\) 进入估计常数，不能在方向趋于共面时忽略。

\ProseParagraphBreak
对于参数曲面 \(\Sigma_j:U_j\to\mathbb R^d\)，也可用切向偏导的楔积 \(Y_j\) 代替单位法向。此时 determinant 下界同时包含参数测度与面测度的 Jacobi 因子，陈述必须说明采用哪一种归一化。

\subsection{几何非退化性}
\label{b2:subsec:250202}

下面采用\term[loomis-whitney-budengshi]{Loomis--Whitney 不等式}[Loomis--Whitney inequality]的下述版本。

\begin{theorem}[Loomis--Whitney 不等式]\label{b2:thm:25-loomis-whitney}
设 \(\pi_j:\mathbb R^d\to\mathbb R^{d-1}\) 删除第 \(j\) 个坐标，\(F_j\ge0\) 可积。则
\[
 \int_{\mathbb R^d}\prod_{j=1}^dF_j(\pi_jx)^{1/(d-1)}\,\mathrm dx
 \le\prod_{j=1}^d\left(\int_{\mathbb R^{d-1}}F_j\right)^{1/(d-1)}.
\]
\end{theorem}
\begin{proof}
对维数归纳。\(d=2\) 时恰好是 Fubini 乘积分离。
设结论在 \(d-1\) 维成立。先对 \(x_d\) 积分，对其余 \(d-1\) 个因子用 Hölder，令
\(G_j=\int F_j\,\mathrm dx_d\)、\(j<d\)；它们各依赖 \(x'=(x_1,\ldots,x_{d-1})\) 中除第 \(j\) 个以外的坐标。于是左侧至多
\[
 \int_{\mathbb R^{d-1}}F_d(x')^{1/(d-1)}
                     \prod_{j<d}G_j(\pi'_jx')^{1/(d-1)}\,\mathrm dx'.
\]
对这两个部分再用指数 \(d-1\) 与 \((d-1)/(d-2)\) 的 Hölder，得到
\[
 \left(\int F_d\right)^{1/(d-1)}
 \left[\int\prod_{j<d}G_j(\pi'_jx')^{1/(d-2)}\,\mathrm dx'\right]^{(d-2)/(d-1)}.
\]
归纳假设与 \(\int G_j=\int F_j\) 给出结论。截断与单调收敛保证所有非负积分的处理合法。
\end{proof}
这条不等式并非普通 Hölder 的同域版本，因为每个因子在一个方向上恒定，其全空间积分会无限大。删除坐标的互补几何是有界性的来源。

\subsection{多线性增益}
\label{b2:subsec:250203}

\begin{proposition}[平直横截延拓]\label{b2:prop:25-flat-extension}
设 \(H_j=n_j^\perp\)、\(n_j\) 为单位向量，矩阵
\(N=(n_1,\ldots,n_d)\) 可逆。以 \(H_j\) 上的欧氏面测度定义延拓，则
\[
 \left\|\prod_{j=1}^dE_jg_j\right\|_{L^{2/(d-1)}(\mathbb R^d)}
 \le|\det N|^{-1/2}\prod_j\|g_j\|_{L^2(H_j)}.
\]
\end{proposition}
\begin{proof}
令 \(x=Ny\)。由于 \(E_jg_j\) 沿 \(n_j\) 方向不变，它在 \(y\) 坐标中不依赖 \(y_j\)。
记其平方模为 \(F_j(\pi_jy)\)，应用 Loomis–Whitney。
在截面 \(y_j=0\) 上，向 \(H_j\) 正交投影的 Jacobi 因子是 \(|\det N|\)：把其他列投影到 \(H_j\)，再以单位列 \(n_j\) 补齐 determinant 即可。
因此由 \(H_j\) 上的 Plancherel，
\(\int F_j=|\det N|^{-1}\|g_j\|_2^2\)。
空间换元贡献 \(|\det N|\)，合并得
\[
 \int\prod_j|E_jg_j|^{2/(d-1)}
 \le|\det N|^{-1/(d-1)}
             \prod_j\|g_j\|_2^{2/(d-1)}.
\]
取 \((d-1)/2\) 次幂即得。当输出指数小于1时按该积分定义准范数，证明仍成立。
\end{proof}
平面没有曲率，却在完整横截乘积中有这个界。线性限制中的平面失败因此不能用来否定多线性横截估计。二者保留的几何信息不同。

% !TeX root = ../../Book2.tex
\section{双线性振荡积分}
\label{b2:sec:2503}

二维曲线的双线性估计可以用一个可逆的频率和映射完整计算。这给出横截增益的直接模型，并显示常数如何随方向靠近而恶化。

\subsection{双线性振荡积分}
\label{b2:subsec:250301}

令 \(\gamma(t)=(t,t^2)\)，
\[
 E_If(x,y)=\int_I e^{2\pi\mathrm i(xt+yt^2)}f(t)\,\mathrm dt.
\]
对两个不交紧区间 \(I,J\)，乘积的 Fourier 侧对应映射
\[
 \Psi(t,s)=(t+s,t^2+s^2),\qquad
 |\det D\Psi|=2|s-t|.
\]
由于 \(t\in I\)、\(s\in J\) 的顺序固定，和与平方和确定无序对 \(\{t,s\}\)，顺序又确定有序对，所以映射在该矩形内一对一。
这项一对一性不能只由 Jacobi 行列式非零推出；后者仅保证局部可逆。

\subsection{双线性估计}
\label{b2:subsec:250302}

\begin{theorem}[分离抛物线片的双线性估计]\label{b2:thm:25-parabola-bilinear}
若 \(\operatorname{dist}(I,J)\ge\rho>0\)，则
\[
 \|E_If\,E_Jg\|_{L^2(\mathbb R^2)}^2
 =\int_{I\times J}\frac{|f(t)|^2|g(s)|^2}{2|t-s|}\,\mathrm dt\,\mathrm ds
 \le(2\rho)^{-1}\|f\|_2^2\|g\|_2^2.
\]
\end{theorem}
\begin{proof}
先取光滑紧支撑输入。换元 \(u=\Psi(t,s)\) 后，乘积为某个密度 \(h(u)\) 的逆 Fourier 变换，且
\(h(\Psi(t,s))=f(t)g(s)/(2|t-s|)\)。
Plancherel 与换元给出
\(\|h\|_2^2=\int |f|^2|g|^2/(2|t-s|)\)，这就是等式。
右侧由分离下界控制。一般 \(L^2\) 输入以光滑函数逼近，两项延拓乘积在 \(L^2\) 中 Cauchy；同时每个延拓在点态积分意义下连续依赖于 \(L^1\) 输入，有限区间上的 \(L^2\) 逼近也给出 \(L^1\) 逼近，故极限确为目标乘积。
\end{proof}
这个定理估计乘积而非单个函数。它没有断言 \(E_I:L^2(I)\to L^4(\mathbb R^2)\) 在允许片自相互作用时成立；对角线 \(t=s\) 的 Jacobi 行列式退化正是该推断的障碍。

\subsection{曲率相互作用}
\label{b2:subsec:250303}

对一般曲线 \(\gamma\)，和映射 Jacobi 行列式为
\(\det(\gamma'(t),\gamma'(s))\)。在足够小且方向分离的两片上，若此 determinant 有正下界、映射具有受控有限重数，同一换元与 Plancherel 给出对应估计，常数还包含重数。
对于抛物线，曲率使不同参数的切向方向改变；对于两条本来就横截的直线，则不需要曲率。

\ProseParagraphBreak
因此双线性理论中的曲率与横截性各有角色。横截性直接保证频率和映射的 Jacobi 行列式，曲率决定同一曲面的不同片能否形成横截方向。高维中两个超曲面的和映射通常有正维纤维，需要沿纤维积分，不能照搬二维的一对一公式。

% !TeX root = ../../Book2.tex
\section{多线性 Fourier 限制}
\label{b2:sec:2504}

完整 \(d\) 线性横截问题是 Loomis–Whitney 模型的振荡推广。平直情形已经证明；一般曲面允许法向变化，需要尺度归纳和多线性管束理论。

\subsection{多线性延拓估计}
\label{b2:subsec:250401}

下面采用\term[bennett-carbery-tao-duoxianxing-xianzhi]{Bennett--Carbery--Tao 多线性限制定理}[Bennett--Carbery--Tao multilinear restriction theorem]的下述版本。

\begin{theorem}[深层结构定理：Bennett--Carbery--Tao 局部多线性限制]\label{b2:thm:25-bct}
设 \(d\geq2\)，\(\Sigma_j:U_j\subset\mathbb R^{d-1}\to\mathbb R^d\) 为固定紧参数片上的光滑参数化，\(C^2\) 范数有统一上界，切向楔积 \(Y_j\) 满足
\(|\det(Y_1(u_1),\ldots,Y_d(u_d))|\ge\nu>0\)。
令 \(E_jg_j(x)=\int_{U_j}e^{2\pi\mathrm ix\cdot\Sigma_j(u)}g_j(u)\,\mathrm du\)。
对每个 \(\epsilon>0\)、\(R\ge1\)，有
\[
 \left\|\prod_{j=1}^dE_jg_j\right\|_{L^{2/(d-1)}(B_R)}
       \le C_\epsilon R^\epsilon\prod_j\|g_j\|_2.
\]
\end{theorem}
\begin{proofsketch}
波包分解将延拓转化为横截管束的组合；尺度归纳把限制界归约到多线性 Kakeya 控制，再以 Gaussian 族的单调性处理管束。误差和尺度迭代产生允许的 \(R^\epsilon\) 损失。
完整管束与归纳证明在此略去。上述版本是
\cite[Theorem 1.16, p.~267; Sections 2--4]{Bennett2006Multilinear}
的 \(L^2\) 端点特例；正式期刊版的假设、指数及尺度损失已核验。
\end{proofsketch}
保留局部球和 \(R^\epsilon\) 很重要。本章没有证明去掉这个因子的全空间端点界。平直模型能无损成立，不等于一般曲面也可直接无损；同样不能把小于1的输出准范数当作 Banach 范数随意使用三角不等式。

\subsection{多线性限制}
\label{b2:subsec:250402}

多线性限制通常用延拓乘积陈述。它与单个线性算子的简单伴随关系不同；固定其他输入后可以得到一个线性配对，再求其伴随，但完整形式涉及多个曲面。
参数测度 \(\mathrm du\) 与面测度的关系为
\(\mathrm d\sigma_j=|Y_j(u)|\,\mathrm du\)。
若参数化的 Jacobi 因子有正上下界，转换输入范数只改变常数；若没有这个下界，不能省略该因子。

\ProseParagraphBreak
最小维数 \(d=2\) 时乘积输出为 \(L^2\)，上一节的曲线和映射证明已给出局部横截模型。高维输出指数 \(2/(d-1)\) 越来越小，反映的是多个独立方向同时使乘积集中区域缩小，而非单个因子的可积性无限改善。

\subsection{对线性限制的反馈}
\label{b2:subsec:250403}

线性输出在一个点很大时，可能主要来自一个帽，或来自多个法向分离的帽。前一种情况可在更小尺度重新规范化，后一种情况可以调用多线性估计。这种“分离方向部分加集中部分”的组织让横截界反馈到线性限制。
实际归约需证明每个点的帽分配、范数求和与尺度损失的递推；仅知道多线性定理不自动给出线性猜想全部范围。

\ProseParagraphBreak
横截 determinant 接近零的帽族不能硬塞进 \(d\) 线性定理。它们可能集中在较低维方向簇中，需要双线性、低维或切向方法。分辨哪些输入属于横截部分，是现代归纳与分割的关键步骤。

% !TeX root = ../../Book2.tex
\section{多线性限制的几何与尺度方法}
\label{b2:sec:2505}

前面的横截模型解释了不同方向怎样改善乘积估计。把这种增益用于许多帽和许多空间区域，还需要尺度递推、波包分解以及对几何组合的控制。本节先证明其中可直接计算的步骤，再陈述抛物面解耦这一深层输入。递推式的证明、单帽展开和多帽组合是不同的任务，应分别检查它们需要的条件。

\subsection{尺度归纳}
\label{b2:subsec:250501}

设 \(A(R)\) 表示半径 \(R\) 球上的某个规范化最佳常数；若定理允许任意球心或一族曲面，便应先对这些参数取上确界。尺度归纳希望把 \(A(R)\) 与较小尺度的常数联系起来。一种常见模型是
\[
 A(R)\leq C_\eta R^\eta A(R^{1/2}).
\]
右侧不只有幂次损失：每迭代一次，还会多出一个 \(C_\eta\)。下面把两者都算出。

\Needspace{6\baselineskip}
\begin{lemma}[平方根尺度递推]\label{b2:lem:25-scale-iteration}
设 \(A:[1,\infty)\to[0,\infty)\) 为有限值函数，满足
\[
 M\coloneqq\sup_{1\leq R\leq2}A(R)<\infty.
\]
假设对每个 \(\eta>0\)，存在与 \(R\) 无关的有限常数 \(C_\eta\geq1\)，使
\[
 A(R)\leq C_\eta R^\eta A(R^{1/2})\quad(R>2).
\]
则对每个 \(\varepsilon>0\)，存在有限常数 \(K_\varepsilon\)，使
\[
 A(R)\leq K_\varepsilon R^\varepsilon\quad(R\geq1).
\]
常数 \(K_\varepsilon\) 可以依赖 \(M\) 及相应的 \(C_\eta\)，不要求它在 \(\varepsilon\downarrow0\) 时保持有界。
\end{lemma}
\begin{proof}
固定 \(\eta>0\) 及 \(R>2\)，令 \(m\) 为满足 \(R^{2^{-m}}\leq2\) 的最小正整数。前 \(m\) 步均可应用递推，故
\[
 \begin{aligned}
 A(R)
 &\leq C_\eta^m
       R^{\eta(1+2^{-1}+\cdots+2^{-(m-1)})}
       A(R^{2^{-m}})\\
 &\leq M C_\eta^m R^{2\eta}.
 \end{aligned}
\]
另一方面，
\[
 m\leq1+\log_2\!\left(\frac{\log R}{\log2}\right).
\]
记 \(b_\eta=\log_2 C_\eta\geq0\)，则
\[
 C_\eta^m
 \leq C_\eta\left(\frac{\log R}{\log2}\right)^{b_\eta}.
\]
所以固定 \(\eta\) 后，全部常数累积只是一个固定的对数幂。

现在给定 \(\varepsilon>0\)，先取 \(\eta=\varepsilon/4\)。这样 \(C_\eta\) 和 \(b_\eta\) 已经固定。函数
\((\log R)^{b_\eta}R^{-\varepsilon/2}\)
在 \(R\geq2\) 上有界：令 \(t=\log R\)，它就是 \(t^{b_\eta}e^{-\varepsilon t/2}\)，指数衰减控制固定幂次。因此
\[
 A(R)\leq K_\varepsilon
          R^{2\eta+\varepsilon/2}
       =K_\varepsilon R^\varepsilon.
\]
初始区间 \(1\leq R\leq2\) 由 \(M\) 控制，增大常数即可。
\end{proof}

顺序在这里很重要：先选 \(\eta\)，再吸收由 \(C_\eta\) 决定的对数幂。若让 \(\eta\) 随 \(R\) 改变，却不知道 \(C_\eta\) 怎样增长，上述吸收便没有统一常数。这个引理也不负责证明具体限制问题所需的递推式；它只说明一旦递推和初始尺度界成立，怎样正确结束迭代。

\ProseParagraphBreak
每一步还必须留在同一允许的曲面类中，包括参数化的导数界、支集大小和横截下界。\cref{b2:prop:25-parabolic-rescaling}给出抛物面的精确重标度；一般曲面的小片重标度则要控制偏离二次模型的误差。如果归纳后曲率或横截常数不断恶化，就不能把下一步仍写成同一个 \(A(R^{1/2})\)，更不能只凭幂次几何级数断言得到任意小损失。

\subsection{波包}
\label{b2:subsec:250502}

在空间半径 \(R\) 上，一个宽度 \(R^{-1/2}\) 的频率帽对应宽约 \(R^{1/2}\)、长约 \(R\) 的管。这与\cref{b2:fig:23-cap-tube}及\cref{b2:prop:23-parabolic-packet}的尺度一致，但单个光滑帽函数还不代表一般 \(L^2\) 输入。下面用帽内 Fourier 级数把一般输入拆成位置不同的波包，并保留系数的平方能量。

\ProseParagraphBreak
为使归一化可直接检查，这一构造使用抛物面的参数测度；面积测度的光滑密度在固定参数片上可以吸收到输入中，但相应的范数比较应随之注明。

\Needspace{6\baselineskip}
\begin{proposition}[单帽的波包展开]\label{b2:prop:25-single-cap-packets}
设整数 \(n\geq1\)，\(R\geq1\)、\(\delta=R^{-1/2}\)，且 \(\xi_0\in\mathbb R^n\)、\(|\xi_0|\leq1\)。取
\[
 g\in L^2(\mathbb R^n),\quad
 \operatorname{supp}g\subset
       \xi_0+\delta[-1/2,1/2]^n,
\]
并以参数测度定义
\[
 E_{\mathrm{par}}g(x',t)
 =\int_{\mathbb R^n}
       e^{2\pi\mathrm i(x'\cdot\xi+t|\xi|^2)}
       \cdot g(\xi)\,\mathrm d\xi.
\]
固定 \(\rho\in C_c^\infty((-1,1)^n)\)，在 \([-1/2,1/2]^n\) 上恒为一。则存在系数 \(c_k\)、\(k\in\mathbb Z^n\)，及只依赖 \(\xi_0,\delta,\rho,k\) 的频率函数 \(\psi_k\)，使
\[
 g=\sum_{k\in\mathbb Z^n}c_k\psi_k
       \quad\text{于 }L^2(\mathbb R^n),\quad
 \sum_k|c_k|^2=\|g\|_2^2.
\]
若 \(W_k=E_{\mathrm{par}}\psi_k\)，则按有限立方格点部分和取极限时，
\[
 E_{\mathrm{par}}g=\sum_kc_kW_k
\]
在整个 \(\mathbb R^{n+1}\) 上一致收敛。对每个 \(N>0\)，在 \(|t|\leq R\) 内还有
\[
 |W_k(x',t)|
 \leq C_{N,n,\rho}\delta^{n/2}
  \left(1+\left|\delta(x'+2t\xi_0)+k/2\right|\right)^{-N}.
\]
可选一个只依赖维数的固定 \(A_0\)，使对任意 \(L>0\)，
\[
 \sup_{|(x',t)|\leq R}
 \left|E_{\mathrm{par}}g(x',t)
       -\sum_{|k|\leq A_0\sqrt R}c_kW_k(x',t)\right|
 \leq C_{L,n,\rho}R^{-L}\|g\|_2.
\]
所有常数均与 \(R,\xi_0,g\) 无关。
\end{proposition}
\begin{proof}
将帽拉回单位盒，令 \(h(u)=g(\xi_0+\delta u)\)。它在
\(Q=[-1,1]^n\) 上支于较小的盒 \([-1/2,1/2]^n\)。正交基
\[
 e_k(u)=2^{-n/2}e^{\pi\mathrm i k\cdot u},\quad k\in\mathbb Z^n,
\]
给出 Fourier 级数
\[
 h=\sum_ka_ke_k\quad\text{于 }L^2(Q),\quad
 \sum_k|a_k|^2=\|h\|_{L^2(Q)}^2
             =\delta^{-n}\|g\|_2^2.
\]
令
\[
 c_k=\delta^{n/2}a_k,\quad
 \psi_k(\xi)=2^{-n/2}\delta^{-n/2}
       \rho\!\left(\frac{\xi-\xi_0}{\delta}\right)
       \cdot e^{\pi\mathrm i k\cdot(\xi-\xi_0)/\delta}.
\]
由于 \(\rho\cdot h=h\)，把 Fourier 级数乘以有界函数 \(\rho\)，再作换元，就得到所述 \(L^2\) 重构与精确系数能量。每个频率函数支于放大一倍的帽，且
\[
 \|\psi_k\|_2=2^{-n/2}\|\rho\|_2.
\]
它们有固定的 \(L^2\) 大小，但并不因这个归一化就彼此正交；系数能量等式来自乘 \(\rho\) 之前的正交 Fourier 基。

换元 \(\xi=\xi_0+\delta u\) 给出原波包的具体公式
\[
 \begin{aligned}
 W_k(x',t)
 ={}&2^{-n/2}\delta^{n/2}
       e^{2\pi\mathrm i(x'\cdot\xi_0+t|\xi_0|^2)}\\
 &{}\cdot
 \int_{\mathbb R^n}\rho(u)\cdot
       e^{2\pi\mathrm i[
          \delta(x'+2t\xi_0)\cdot u
          +t\delta^2|u|^2+k\cdot u/2]}\,\mathrm du.
 \end{aligned}
\]
频率部分和的误差支于固定的放大帽，其 \(L^2\) 范数趋零，故其 \(L^1\) 范数也趋零。延拓积分核的模为一，因此输出误差的一致范数至多为这个 \(L^1\) 误差。这证明全空间的一致收敛，并没有假定一般 \(L^2\) 函数的 Fourier 级数逐点收敛。

若 \(|t|\leq R\)，则 \(|t\delta^2|\leq1\)。振幅
\(\rho(u)\cdot e^{2\pi\mathrm i t\delta^2|u|^2}\)
的支集固定，各固定阶导数一致有界。对剩余的线性频率
\[
 z+k/2,\quad z=\delta(x'+2t\xi_0),
\]
反复分部积分，即得所述任意阶衰减。

还需把这些尾部实际求和。对任意 \(z\in\mathbb R^n\)，格点间距固定，半径 \(2^j\) 的球内至多有 \(C_n2^{jn}\) 个点。因此当 \(2N>n\) 时，
\[
 \sup_z\sum_{k\in\mathbb Z^n}
       (1+|z+k/2|)^{-2N}<\infty.
\]
离散 Cauchy--Schwarz 与系数能量给出
\[
 \sum_k|c_k|\cdot|W_k(x',t)|
 \leq C_{N,n,\rho}\delta^{n/2}\|g\|_2
 \quad(|t|\leq R).
\]
所以该时间片内的波包和还逐点绝对收敛。

在球 \(|(x',t)|\leq R\) 内，因 \(|\xi_0|\leq1\)，有
\(|z|\leq3\sqrt R\)。取固定 \(A_0\geq12\)，若 \(|k|>A_0\sqrt R\)，则
\[
 |z+k/2|\geq|k|/4.
\]
同样按格点环带求和，得到
\[
 \begin{aligned}
 \sum_{|k|>A_0\sqrt R}|c_k|\cdot|W_k(x',t)|
 &\leq C_{N,n,\rho}\delta^{n/2}\|g\|_2
        R^{n/4-N/2}\\
 &=C_{N,n,\rho}R^{-N/2}\|g\|_2.
 \end{aligned}
\]
取 \(N>\max\{n/2,2L\}\)，便得到命题的局部误差界。
\end{proof}

在这个构造中，主要管区域为
\[
 |x'+2t\xi_0+k/(2\delta)|\lesssim\delta^{-1},\quad
 |t|\leq R.
\]
其宽度约为 \(R^{1/2}\)，长度约为 \(R\)，长轴沿抛物面的法向 \((-2\xi_0,1)\)。帽的中心决定方向，格点 \(k\) 决定位置；同一帽中的管彼此平行。空间球内只需保留 \(O(R^{n/2})\) 个位置指标，其余项可以按任意预先指定的负幂控制。

\ProseParagraphBreak
这里有三个不同的结论。频率重构在 \(L^2\) 中成立；紧频率支集又使延拓部分和在全空间一致收敛；管外快速衰减及上述有限格点截断则限于 \(|t|\leq R\) 或 \(|(x',t)|\leq R\)。最后一项不能无条件推广到任意大时间。波包并非管的示性函数：公式中的相位与系数仍须保留，截到粗管时产生的尾部也须按需要的阶数求和。

\ProseParagraphBreak
这只是单张帽的展开。若把许多帽相加，各帽之间有不同的法向，不能从一张帽的系数能量等式推出 BCT 或解耦。为了估计这些组合，接下来需要同时看管与空间分区的关系，以及不同帽输出的范数怎样相加。

\subsection{多项式分割}
\label{b2:subsec:250503}

多项式分割选择一个多项式，把大部分空间分成若干胞腔，并将零集的邻域单独保留。它希望让各胞腔承担可控的质量，同时利用管的穿越次数限制重复计算。选择分割多项式及控制各胞腔质量属于额外的存在性结论，不能由直线的根数直接推出。

\ProseParagraphBreak
\cref{b2:lem:23-polynomial-line-crossing}证明：次数至多 \(D\) 的非零多项式，其零集之外的胞腔，至多有 \(D+1\) 个能被同一条不全含于零集的直线经过。对于半径为 \(\rho\) 的粗管，要把胞腔缩到与零集距离大于 \(\rho\) 的部分，才能把管的相遇转回中心直线的穿越。原因是管中这样的点到中心线的连接线段避开零集，故其中心线点仍在同一胞腔。未经缩减的胞腔与零集邻域，不由这个次数界控制。

即使多项式只有一次，零集邻域中的管也会表现出两种不同情形。平面模型足以把角度和停留长度算清。

\Needspace{6\baselineskip}
\begin{proposition}[管与平面厚邻域]\label{b2:prop:25-plane-wall}
设整数 \(d\geq2\)，\(\nu,v\in S^{d-1}\)，\(a\in\mathbb R^d\)、\(b\in\mathbb R\)，且 \(L,r,\rho>0\)。令
\[
 H=\{\,x:x\cdot\nu=b\,\},\quad
 W_r(H)=\{\,x:|x\cdot\nu-b|\leq r\,\},
\]
\[
 \gamma(s)=a+sv\quad(0\leq s\leq L),\quad
 T_\rho=\bigcup_{0\leq s\leq L}\overline{B(\gamma(s),\rho)}.
\]
记可能与平面厚邻域相遇的轴参数集合为
\[
 I=\{\,s\in[0,L]:
       \overline{B(\gamma(s),\rho)}\cap W_r(H)\neq\varnothing\,\}.
\]
以一维 Lebesgue 测度记 \(|I|\)。若 \(v\cdot\nu\neq0\)，则
\[
 |I|\leq\min\left\{L,\frac{2(r+\rho)}{|v\cdot\nu|}\right\}.
\]
若 \(v\cdot\nu=0\)，则 \(I=[0,L]\) 或 \(I=\varnothing\)，前一种情形恰好对应 \(|a\cdot\nu-b|\leq r+\rho\)。特别地，若轴线在 \(H\) 内且 \(\rho\leq r\)，则整条管 \(T_\rho\) 都包含于 \(W_r(H)\)。
\end{proposition}
\begin{proof}
球 \(\overline{B(\gamma(s),\rho)}\) 与平面厚邻域相遇，当且仅当其球心到平面的距离至多为 \(r+\rho\)。于是
\[
 I=\{\,s\in[0,L]:
       |a\cdot\nu-b+s(v\cdot\nu)|\leq r+\rho\,\}.
\]
若斜率 \(v\cdot\nu\neq0\)，不等式在整个实轴上给出长度
\(2(r+\rho)/|v\cdot\nu|\) 的区间，与 \([0,L]\) 相交后得到所述上界。若斜率为零，条件与 \(s\) 无关，所以参数集合全有或全无。轴线位于平面内时，管中每点到该平面的距离不超过 \(\rho\)，最后一项随即成立。
\end{proof}

若 \(\vartheta\) 是管轴与平面的夹角，则 \(|v\cdot\nu|=|\sin\vartheta|\)。取波包尺度 \(L=R\)、\(r\) 与 \(\rho\) 都约为 \(R^{1/2}\)，固定正角度给出的停留长度只有 \(O(R^{1/2})\)；当角度小到 \(R^{-1/2}\) 的量级时，上界可以达到整条管的长度 \(R\)。平行而贴墙的管更可能一直留在墙内，不能把它作为普通横穿管处理。

\ProseParagraphBreak
这说明零集邻域中的分析不能只问“中心线与零集交了几次”。横穿部分可利用角度下界缩短停留区域；近切向部分则可能沿低维几何聚集，需要不同的估计。对高次零集，局部切平面会随位置改变，还会有奇点及多个局部片，平面模型本身不处理这些困难。

\ProseParagraphBreak
Guth 的《A Restriction Estimate Using Polynomial Partitioning》\cite{Guth2016Restriction}给出了这种方法的正式限制应用。多项式次数一方面增加胞腔数，另一方面也增加穿越预算和零集邻域的复杂度，必须放进尺度递推中选择。本节没有重建分割定理、切向管分析或完整限制范围证明，因而不从上述次数与角度计算推出新的线性限制估计。

\subsection{解耦}
\label{b2:subsec:250504}

单帽波包保留帽内的相干性；解耦则控制许多帽叠加后的输出。下面采用\term[jieou-dingli]{解耦定理}[decoupling theorem]的抛物面临界版本。帽的尺度、分割函数和空间权都属于陈述的一部分。

\Needspace{6\baselineskip}
\begin{theorem}[深层结构定理：抛物面临界解耦]\label{b2:thm:25-decoupling}
设整数 \(d\geq2\)、\(R\geq1\)，令
\[
 S=\{\,(\xi',|\xi'|^2):|\xi'|\leq1\,\}\subset\mathbb R^d,\quad
 p_c=\frac{2(d+1)}{d-1},\quad \delta=R^{-1/2}.
\]
曲面取面积测度 \(\sigma\)，并定义
\[
 Eg(x)=\int_Se^{2\pi\mathrm i x\cdot\xi}
       \cdot g(\xi)\,\mathrm d\sigma(\xi)
 \quad(g\in L^2(S,\sigma)).
\]
在参数域上取尺度 \(\delta\) 的规则方格帽覆盖，并固定非负光滑单位分解 \((\chi_\theta)_\theta\)：每个 \(\chi_\theta\) 支于相应方格的固定倍数放大，其支集重叠次数一致有界，
\[
 \sum_\theta\chi_\theta=1\quad\text{于 }S,\quad
 |\partial^\alpha\chi_\theta|\leq C_\alpha\delta^{-|\alpha|}.
\]
导数在参数坐标中计算，\(C_\alpha\) 与 \(R,\theta\) 无关。记
\[
 E_\theta g\coloneqq E(\chi_\theta\cdot g).
\]
若 \(B_R=B(c,R)\subset\mathbb R^d\)，令
\[
 w_{B_R}(x)=\left(1+\frac{|x-c|}{R}\right)^{-10d},\quad
 \|F\|_{L^{p_c}(w_{B_R})}
 =\left(\int_{\mathbb R^d}|F(x)|^{p_c}
          w_{B_R}(x)\,\mathrm dx\right)^{1/p_c}.
\]
则对每个 \(\varepsilon>0\)，存在 \(C_\varepsilon<\infty\)，使
\[
 \|Eg\|_{L^{p_c}(B_R)}
 \leq C_\varepsilon R^\varepsilon
       \left(\sum_\theta
          \|E_\theta g\|_{L^{p_c}(w_{B_R})}^2
       \right)^{1/2}.
\]
常数可以依赖 \(d,\varepsilon\) 及固定分割的重叠、导数界，与 \(R,c,g\) 无关。
\end{theorem}
\begin{proofsketch}
频率薄邻域中的全空间解耦经空间局部化得到上述延拓形式；固定的光滑帽分割用于定义各项，有限重叠只改变常数。完整证明结合多线性限制、抛物重标度和反复尺度迭代，最后保留任意小幂次损失。这里略去具体的多尺度递推与闭合，不能仅凭\cref{b2:lem:25-scale-iteration}替代那些步骤。

正式来源为 Bourgain 与 Demeter 的《The Proof of the \(\ell^2\) Decoupling Conjecture》%
\cite[Theorem 1.1; equation (7), pp.~352--355]{BourgainDemeter2015Proof}。
该版本要求正定第二基本形式，临界指数如上。标准抛物面满足这些条件；不能把同一结论未经修改套到第二基本形式不定或退化的曲面。
\end{proofsketch}

定义中的 \(E_\theta g\) 始终使用 \(\chi_\theta\cdot g\)，所以
\(Eg=\sum_\theta E_\theta g\) 精确成立。若改用互不相交的可测帽示性函数，应从头采用那种分割的版本；不能既把帽当作重叠覆盖，又让它们的示性函数相加充当一。权 \(w_{B_R}\) 在球内有统一正下界，在球外保留快速衰减尾部，这与波包没有严格管支集相适应。换成双边可比的权只改变常数，不能任意删去球外部分而保持同一陈述。

\ProseParagraphBreak
还应区分两个容易混写的右端。对有限个函数 \(F_\theta\) 及任意正测度 \(\mu\)，当 \(2\leq p<\infty\) 且各项属于 \(L^p(\mu)\) 时，Minkowski 不等式给出
\[
 \begin{aligned}
 \left\|\left(\sum_\theta|F_\theta|^2\right)^{1/2}\right\|_{L^p(\mu)}^2
 &=\left\|\sum_\theta|F_\theta|^2\right\|_{L^{p/2}(\mu)}\\
 &\leq\sum_\theta
       \bigl\||F_\theta|^2\bigr\|_{L^{p/2}(\mu)}
 =\sum_\theta\|F_\theta\|_{L^p(\mu)}^2.
 \end{aligned}
\]
因此点态平方函数的范数至多为各函数范数的 \(\ell^2\) 和。把解耦定理右端换成前者，会提出一个通常更强的估计，不能由显示的解耦不等式直接推出。

\ProseParagraphBreak
一个有限模型可把差别算出。取 \(\mathbb R^d\) 中 \(N\) 个互不相交、体积为一的方块 \(Q_1,\ldots,Q_N\)，令 \(F_j=\mathbf1_{Q_j}\)。则
\[
 \left\|\left(\sum_{j=1}^N|F_j|^2\right)^{1/2}\right\|_p=N^{1/p},
 \quad
 \left(\sum_{j=1}^N\|F_j\|_p^2\right)^{1/2}=N^{1/2}.
\]
在临界 \(p=p_c\) 处，两者的比值为 \(N^{1/(d+1)}\)。这个例子只比较两种范数结构，没有把方块示性函数宣称为曲面帽的延拓。若各个 \(F_j\) 完全相同，两者又相等；差别取决于各项在空间中的分布，不能用一句“正交”概括。

\ProseParagraphBreak
在 \(p=2\) 时，上述两种范数由积分的线性性相等；非平方指数处则存在实际差距。抛物面临界解耦的内容是以 \(R^\varepsilon\) 控制许多频率方向的叠加，并保留加权球和指定帽尺度。单帽的系数能量、平面穿越计算及尺度递推引理各解决其中一部分问题，不能单独代替这项多帽估计。


\begin{chaptersummary}
多线性非退化取决于相位与投影族，不能仅看标量 Hesse 矩阵。完整横截的平直模型由 Loomis–Whitney 控制，二维曲线模型由和映射 Jacobi 行列式控制。一般曲面需要波包和尺度归纳；解耦进一步控制各帽输出的平方和。局部半径、横截下界与任意小幂次损失均是正式结论的一部分。
\end{chaptersummary}

\BookExercises

% 自拟：围绕本章已经建立的工具，改变条件或组合方法。
\begin{optionalexercise}\label{b2:ex:25-new-1}
对投影 \(\pi_1(x,y)=x,\pi_2(x,y)=y\)，证明相位 \(x^2+y^2\) 不可能提供统一的正幂次振荡衰减。
\end{optionalexercise}
\begin{solution}
选固定光滑 \(g,h\) 使非振荡积分非零，令输入分别乘 \(e^{-2\pi\mathrm i\lambda x^2}\)、\(e^{-2\pi\mathrm i\lambda y^2}\)。输入范数不变，相位完全抵消，输出保持一个非零常数，排除任何趋零的统一界。
\end{solution}

% 自拟：围绕本章已经建立的工具，改变条件或组合方法。
\begin{optionalexercise}\label{b2:ex:25-new-2}
在 \(d=3\) 中写出 Loomis–Whitney 不等式，并用两次 Cauchy–Schwarz 证明它。
\end{optionalexercise}
\begin{solution}
不等式为 \(\int[F_1(y,z)F_2(x,z)F_3(x,y)]^{1/2}\le\prod(\int F_j)^{1/2}\)。先对 \(z\) 用 Cauchy–Schwarz，得到 \(G_1(y)^{1/2}G_2(x)^{1/2}\)；再对 \((x,y)\) 把它与 \(F_3^{1/2}\) 用 Cauchy–Schwarz，余项积分为 \((\int G_1)(\int G_2)\)，Fubini 给出结论。
\end{solution}

% 自拟：围绕本章已经建立的工具，改变条件或组合方法。
\begin{optionalexercise}\label{b2:ex:25-new-3}
给出三个两两不平行但不满足三维完整横截条件的单位法向。
\end{optionalexercise}
\begin{solution}
取 \(e_1,e_2,(e_1+e_2)/\sqrt2\)。它们两两不平行，但全部位于同一平面，三列 determinant 为0。因此两两角度下界不蕴含三维法向张成条件。
\end{solution}

% 自拟：围绕本章已经建立的工具，改变条件或组合方法。
\begin{optionalexercise}\label{b2:ex:25-new-4}
对 \(I=[0,1]\)、\(J=[2,3]\)，证明抛物线双线性算子范数不超过 \(1/\sqrt2\)。
\end{optionalexercise}
\begin{solution}
两区间距离为1。精确公式中 \(2|t-s|\ge2\)，故乘积平方范数至多 \(\tfrac12\|f\|_2^2\|g\|_2^2\)，取平方根即得。这只是方便上界，没有声称常数被某个 \(L^2\) 输入达到。
\end{solution}

% 自拟：围绕本章已经建立的工具，改变条件或组合方法。
\begin{optionalexercise}\label{b2:ex:25-new-5}
说明抛物线和映射的局部 Jacobi 行列式非零为何还不足以单独写出全局换元等式。
\end{optionalexercise}
\begin{solution}
局部可逆映射可能有多个逆像，换元时须对逆像求和。这里和与平方和只确定无序对，若两输入区间重合，通常出现交换的两个逆像；分离且有固定顺序的区间排除了交换，才得到一对一公式。一般曲线要另行控制重数。
\end{solution}

% 自拟：围绕本章已经建立的工具，改变条件或组合方法。
\begin{optionalexercise}\label{b2:ex:25-new-6}
在 \(d=4\) 时，多线性限制的乘积输出指数是什么？说明怎样在这个指数中正确写有限求和的界。
\end{optionalexercise}
\begin{solution}
指数为 \(2/3\)。它是准范数，不能用普通 Minkowski；正确的基本界为 \(\|\sum_\ell h_\ell\|_{2/3}^{2/3}\le\sum_\ell\|h_\ell\|_{2/3}^{2/3}\)，来自逐点 \(|\sum h_\ell|^{2/3}\le\sum|h_\ell|^{2/3}\) 再积分。
\end{solution}

% 自拟：围绕本章已经建立的工具，改变条件或组合方法。
\begin{optionalexercise}\label{b2:ex:25-new-7}
给出输入频率支撑互不相交、但某些乘积输出频率支撑重合的简单例子。
\end{optionalexercise}
\begin{solution}
在一维取分别位于整数附近的窄频率区间。输入中心 \(0,1,2,3\) 各异，但中心 \(0+3\) 与 \(1+2\) 相同，两个乘积的 Fourier 支撑都在3附近。故输入正交不意味着乘积正交，需核对和集合重叠。
\end{solution}

% 自拟：围绕本章已经建立的工具，改变条件或组合方法。
\begin{optionalexercise}\label{b2:ex:25-new-8}
比较解耦右侧 \((\sum\|u_\theta\|_p^2)^{1/2}\) 与点态平方函数的 \(L^p\) 范数，\(p\ge2\) 时哪一个更大？
\end{optionalexercise}
\begin{solution}
对 \(p/2\ge1\) 使用 Minkowski，得 \(\|\sum|u_\theta|^2\|_{p/2}\le\sum\||u_\theta|^2\|_{p/2}\)。取平方根得到 \(\|(\sum|u_\theta|^2)^{1/2}\|_p\le(\sum\|u_\theta\|_p^2)^{1/2}\)。所以用点态平方函数作右侧是更强的结论，不能由普通解耦直接替换。
\end{solution}
